\documentclass[12pt]{article}
\usepackage{amsmath}
\usepackage{amsthm}
\usepackage{newtx}
\usepackage[margin=1.5in]{geometry}

\DeclareMathOperator{\val}{val}

\newtheorem{conj}{Conjecture}
\newtheorem{prop}[conj]{Proposition}
\newtheorem{thm}[conj]{Theorem}
\newtheorem{corr}[conj]{Corrolary}


\newcommand{\two}{\textsc{two}}

\newcommand{\barv}{\bar{v}}
\newcommand{\barw}{\bar{w}}

\newcommand{\rhov}{\rho v}
\newcommand{\rhow}{\rho w}

\newcommand{\rhobarv}{\rho \barv}
\newcommand{\rhobarw}{\rho \barw}

%\DeclareMathOperator{\weight}{wt}

\newcommand{\weight}[1]{\mathrm{wt}(#1)}

\newcommand{\newthought}{%
  \smallskip\centerline{\rule{2in}{0.5pt}}\smallskip}

\begin{document}

\begin{center}
  \textsc{\large Maximum Chains in Paul's Poset }
\end{center}

\begin{itemize}
\item For a positive integer $b$ let $W_b$ be the set of all
  length-$b$ binary words.

\item For $w \in W_b$, let $\val w$ be the integer represented in
  binary by $w$; that is
  \[
  \val w = \sum_{i=1}^b w_i 2^{b-i} .
  \]
  For example, $\val (1,1,0,1) = 1101_{\textsc{two}} = 13$.

\item For $w \in W_b$, let $\rho(w)$ be the \emph{reversal} of
  $w$. That is $$\rho(d_1,d_2,\ldots,d_b) =
  (d_b,d_{b-1},\ldots,d_1).$$

\item For words $v,w\in W_b$, define $v \prec w$ provided $\val v <
  \val w$ and $\val \rho(v) < \val \rho(w)$. Note that $(W_b,\prec)$
  is a poset.

\item A \emph{chain} is a subset of $W_b$ containing pairwise
  comparable words.



\item Let $c(b)$ be the size of a maximum chain in $W_b$, i.e., the
  height of $W_b$.

\end{itemize}

\begin{thm}\label{thm:main}
  For a positive integer $b$,
  \[
  c(b) = 
  \begin{cases}
    3\cdot 2^{b/2-1} & \text{if $b$ is even and} \\
    2^{(b+1)/2} & \text{if $b$ is odd}.
  \end{cases}
  \]
\end{thm}



\begin{prop}[Chain lower bound]\label{prop:c-lower}
    For a positive integer $b$,
  \[
  c(b) \ge
  \begin{cases}
    3\cdot 2^{b/2-1} & \text{if $b$ is even and} \\
    2^{(b+1)/2} & \text{if $b$ is odd}.
  \end{cases}
  \]
\end{prop}

\begin{proof}
  Observe that if $w_1\prec w_2 \prec \cdots \prec w_n$ is a chain in $W_b$, then 
  $$
  0w_10 \prec 0w_20 \prec \cdots \prec 0w_n0 \prec 1w_11\prec
  1w_21\prec \cdots \prec 1w_n1
  $$ is a chain
  in $W_{b+2}$ that is twice as long.

  We see that $c(1)=2$ because $0\prec1$ is a chain in $W_1$ and that
  $c(2)=3$ because $00\prec01\prec11$ is a chain in $W_2$ (and easy to
  check we cannot make a longer one).

  The result now follows by induction on $b$. 
\end{proof}


\begin{proof}[Alternative proof sketch for odd $b$]
  The palindromes of $W_b$ form a chain. When $b$ is odd, the number
  of palindromes is $2^{(b+1)/2}$. This argument doesn't work for
  even $b$.
\end{proof}

\bigbreak
\noindent\textbf{I and H transformations}

\begin{itemize}
\item A \emph{type I} transformation replaces a $0$ in $v$ with a $1$.

  Example: $001\underline0 101 \to 001 \underline1 101$. 

\item A \emph{type H} transformation replaces a substring in $v$ of
  the form\footnote{The substring begins and ends with a $0$ and has
    one or more consecutive $1$s in between.}  $011\dots10$ with its
  complement, i.e., $100\ldots01$.

  Example: $0\underline{0110}01 \to 0 \underline{1001} 01$.
\end{itemize}

\begin{prop}
  Suppose $v,w\in W_b$ and $w$ is derived from $v$ by an I
  transformation, than $v \prec w$.
\end{prop}


\begin{proof}
  Let $v=(x,0,y)$ and $w=(x,1,y)$. Then $\val w = \val v + 2^k > \val
  v$ for some $k$. Likewise, $\val\rho w = \val\rho v + 2^j > \val\rho
  v$ for some $j$. Therefore $v \prec w$.
\end{proof}

\begin{prop}
  Suppose $v,w\in W_b$ and $w$ is derived from $v$ by an H
  transformation, than $v \prec w$.
\end{prop}

\begin{proof}
  We show that an H-transform increases the numerical value of a
  word. Since H-transforms are symmetrical, it also increases the
  value of the word's reversa

  Suppose that the $011\dots10$ pattern  has its first $0$ in the $2^a$
  column and the last $0$ is in the $2^b$ column (where $a>b$)

  The H-transform increases the value of the word by $2^b+2^a$ and
  decreases its value by $2^{b+1} + 2^{b+2} + \cdots + 2^{a-1}$ which
  equals $2^a - 2^{b+1}$, so the net change is
  \[
  \left(2^a+2^b\right) - \left(2^a - 2^{b+1}\right) = 2^b + 2^{b+1}
  \]
  which is positive. 
\end{proof}



\begin{corr}
  Let $v,w \in W_b$. If $w$ is obtained from $v$ by a series of one or
  more I and H transformations, then $v\prec w$. \qed
\end{corr}

The interesting part is the converse.


\begin{prop}\label{prop:HI}
  Let $v,w \in W_b$ with $v \prec w$. Then $w$ is obtained from $v$
  by a sequence of I and H transformations. 
\end{prop}

\begin{proof}
  The proof is by strong induction on $b$. It is easy to verify for
  $W_1$, $W_2$, and $W_3$. Suppose the result has been shown for all
  $W_i$ with $i<b$. 

  Let $v,w \in W_b$. We consider the first and last bits of both
  words, and break into cases. All told there are 16 possibilities:
  Four possibilities for $v$: $(0,\barv,0)$, $(0,\barv,1)$,
  $(1,\barv,0)$, and $(1,\barv,1)$, and likewise for $w$.


  Cases as follows:

  \begin{enumerate}
  \item $v=(0,\barv,0)$ and $w=(0,\barw,0)$: Induction on
    $\barv\prec\barw$. 

  \item $v=(0,\barv,0)$ and $w=(0,\barw,1)$: Induction on
    $(\barv,0)\prec(\barw,1)$. 


  \item $v=(0,\barv,0)$ and $w=(1,\barw,0)$: Induction on
    $(0,\barv)\prec(1,\barw)$. 


  \item $v=(0,\barv,0)$ and $w=(1,\barw,1)$: Use I transformations to
    go from $v=(0,\barv,0)$ to $(0,11\ldots1,0)$, then one H
    transformation to $(1,00\ldots0,1)$, and then I transformations
    to $w=(1,\barw,1)$.

    \newthought
    
  \item $v=(0,\barv,1)$ and $w=(0,\barw,0)$: Contradiction: $\val\rho v >
    \val\rho w$.


  \item $v=(0,\barv,1)$ and $w=(0,\barw,1)$: Induction on
    $(\barv,1)\prec(\barw,1)$. 


  \item $v=(0,\barv,1)$ and $w=(1,\barw,0)$: Contradiction:
    $\val\rho v > \val\rho w$. 


  \item $v=(0,\barv,1)$ and $w=(1,\barw,1)$: Induction on
    $(0,\barv)\prec(1,\barw)$. 

    
    \newthought

  \item $v=(1,\barv,0)$ and $w=(0,\barw,0)$: Contradiction: $\val v >
    \val w$. 

  \item $v=(1,\barv,0)$ and $w=(0,\barw,1)$:  Contradiction: $\val v >
    \val w$. 

  \item $v=(1,\barv,0)$ and $w=(1,\barw,0)$: Induction on
    $(1,\barv)\prec(1,\barw)$. 

  \item $v=(1,\barv,0)$ and $w=(1,\barw,1)$: Induction on
    $(1,\barv)\prec(1,\barw)$.

    \newthought
    
  \item $v=(1,\barv,1)$ and $w=(0,\barw,0)$: Contradiction: $\val v >
    \val w$.


  \item $v=(1,\barv,1)$ and $w=(0,\barw,1)$: Contradiction: $\val v >
    \val w$.


  \item $v=(1,\barv,1)$ and $w=(1,\barw,0)$: Contradiction: $\val \rho
    v > \val \rho w$.


  \item $v=(1,\barv,1)$ and $w=(1,\barw,1)$: Induction on
    $\barv\prec\barw$. \qedhere

  \end{enumerate}
\end{proof}



Recall that, in a poset, $v$ is \emph{covered by} $w$ provided $v
\prec w$ and there is no element $x$ with $v \prec x \prec w$. 


\begin{corr}
  For $v,w \in W_b$, $v$ is covered by $w$ if and only if $w$ is
  obtained from $v$ by a type I or type H transformation.\qed
\end{corr}

In particular, if $w_1 \prec w_2 \prec \cdots \prec w_t$ is a maximum
chain in $W_b$, then each $w_{i+1}$ is obtained from $w_i$ by a type~I
or a type~H transformation. 


\bigbreak
\noindent\textbf{Weighting function}

For a binary word $v$ we define an indexing system and a weighting
function $\weight v$ as follows.

We index the entries of words from their middle outwards. 
\begin{itemize}
\item If $v$ has odd length ($b=2k+1$) the middle bit is index
  $0$. Bits to the left are indexed $-1,-2,\ldots,-k$ and bits to the
  right are indexed $1,2,\ldots k$.

\item If $v$ has even length ($b=2k$) the bits are index (from left to
  right) as $-k,-(k-1), \ldots,-1,1,2,\ldots,k-1,k$. There is no index
  $0$. 
\end{itemize}

We associate a weight with each index.
\begin{itemize}
\item If $v$ has odd length, the middle three bits have weight
  $1$; these are the entries at indices $-1,0,1$. 

  Moving outwards, the bits at index $\pm j$ (where $j>1$) have weight
  $2^{j-1}$. 

  For example, for a $9$-bit word, the weights (from left to right)
  are $8,4,2,1,1,1,2,4,8$.

\item If $v$ has even length, the middle two bits (at indices $\pm1$)
  have weight $1$. Moving outwards, the bits at index $\pm j$
  (where $j>1$) have weight $\frac32\cdot2^{j-2}$. 

  For example, for an $8$-bit word, the weights (from left to right)
  are $6,3,\frac32,1,1,\frac32,3,6$.
\end{itemize}

We define the weight of a word $v$ to be the sum of the weights
associated with the 1s in $v$. Think of this as the dot product of the
list of weights and the list of digits. 


Examples:
\[
\begin{aligned}
  \weight{110101110} =& [8,4,2,1,1,1,2,4,8]\cdot \\
  &[1,1,0,1,0,1,1,1,0] 
  = 8+4+1+1+1+2+4 = 21 
  \\[6pt]
  \weight{110101} =& [1,1,0,1,0,1] \cdot\\
  &[3,\tfrac32,1,1,\tfrac32,3] = 3 + \tfrac32 + 1 + 3 = 8.5
\end{aligned}
\]

\begin{prop}\label{prop:wt+1}
  If $w$ is obtained from $v$ by an I or H transformation, then 
  $\weight w \ge \weight v + 1$.
\end{prop}

\begin{proof}
  If $w$ is obtained from $v$ by an I transformation, then some bit of
  $v$ is changed from a $0$ to a $1$, and all other bits are
  unchanged. Since all weights are at least $1$, we have 
  $\weight w \ge \weight v + 1$.

  Suppose $w$ is obtained from $v$ by an H transformation. 

  \bigbreak

  \noindent\textbf{Case I:} $b$ is odd. 

  We consider the locations of the ends of the H in $v$, i.e., the
  zeros in the pattern 011\ldots10.

  \begin{itemize}
  \item \emph{The two zeros do not surround the center (index $0$) bit.}
    We assume that the 0s are located at indices $a$ and $b$ with
    $0<a<a+1<b$ (therefore $b-1>a$).


    The addition to the weight is $2^{a-1}+2^{b-1}$. The subtraction
    is $2^a+2^{a+1}+\cdots+2^{b-2}$ which equals
    \[
    2^a\left(1+2+\cdots+2^{b-a-2}\right) = 2^a(2^{b-a-1}-1) = 
    2^{b-1} - 2^a .
    \]
    Thee net change in weight is 
    \[
    \left(2^{a-1}+2^{b-1}\right) -\left(2^{b-1}-2^a\right) =
    2^{a-1}+2^a > 1.
    \]


    \item \emph{One of the zeros is at index $0$ and the other is at
        index $a$}. We may assume $a\ge2$. 

      Addition to the weight is $1 + 2^{a-1}$. Subtraction is
      $1+2+\cdots+2^{a-2}  = 2^{a-1}-1$. 
      
      The net change in weight is 
      \[
      \left(1+2^{a-1}\right) - \left( 2^{a-1}-1 \right) = 2 > 1.
      \]

    \item \emph{One of the zeros is at index $-1$ and the other is at
        index $a\ge1$}. 
      
      Addition to the weight is $1+2^{a-1}$.  Subtraction from the
      weight is 
      \[
      1 + \left(1+2+\cdots+2^{a-2}\right) = 1 + \left(2^{a-1}-1\right)
      = 2^{a-1}.
      \]
      The net change in weight is $1+2^{a-1}-2^{a-1}=1$. 

      \bigbreak
      We have consider the cases where the H configuration does not
      overlap the center, has the center as an end point, or overlaps
      the center with one end immediately next to center. 

    \item \emph{One of the zeros is to the left of center and the
        other is to the right of center, with neither next to the
        center.} We may assume that the zeros are at indices $-a$ and
      $b$, where $a,b\ge2$.

      Addition to the weight is $2^{a-1}+2^{b-1}$. Subtraction from
      the weight is
      \[
      \underbrace{
        \left(1+2+\cdots+2^{a-2}
        \right)
      }_{\text{left of center}} + 1 + 
      \underbrace{
        \left(1+2+\cdots+2^{b-2}
        \right)
      }_{\text{right of center}}
      \]
      which equals 
      \[
      \left(2^{a-1}-1 \right) + 1 + \left(2^{b-1}-1\right) = 2^{a-1}+2^{b-1}-1.
      \]

      The net change in weight is $\left(2^{a-1}+2^{b-1}\right) - 
      \left(2^{a-1}+2^{b-1}\right) = 1$. 
  \end{itemize}

  
  

  \noindent\textbf{Case II:} $b$ is even.

  As before, we consider the location of the 0s in the H configuration
  01\ldots10. 

  \begin{itemize}
  \item \emph{The H configuration does not include either middle bit.}
    That is, the two 0s are at indices $a$ and $b$ with $1<a<a+1<b$. 

    Addition to the weight: $\frac32\cdot 2^{a-2}+ \frac32\cdot
    2^{b-2}=\frac32(2^{a-2}+2^{b-2})$. 

    Subtraction from the weight:
    \[
    \begin{aligned}
      \frac32\left(2^{a-1}+\cdots+2^{b-3}\right) &=
      \frac32\cdot2^{a-1}\left(1+\cdots+2^{b-a-2}\right) \\
      &=
      \frac32\cdot 2^{a-1}\left(2^{b-a-1}-1\right) \\
      &=
      \frac32\left(2^{b-2}-2^{a-1}\right)
    \end{aligned}
    \]

    Net addition to weight is
    \[
    \frac32\left(2^{a-2}+2^{b-2}\right)
    - 
    \frac32\left(2^{b-2}-2^{a-1}\right) = 
    \frac32\left(2^{a-2}+2^{a-1}\right)
    > 1.
    \]

  \item \emph{One of the zeros is a middle bit of the word and the
      other zero is on the same side of the middle.}  

    That is, the zeros are at index 1 and $a$ where $a>2$.

    Addition: $1 + \frac32\cdot 2^{a-2}$. 

    Subtraction:
    \[
    \frac32\left(1 + \cdots + 2^{a-3}\right) 
    =\frac32\left( 2^{a-2}-1 \right)
    \]

    Net change:
    \[
    \left(1 + \frac32\cdot 2^{a-2} \right) - \frac32\left( 2^{a-2}-1
    \right) = \frac52 > 1. 
    \]


  \item \emph{One of the zeros is at index $-1$ and the other is at
      index $a>1$.} 

    Addition: $1 + \frac32\cdot 2^{a-2}$. 

    Subtraction: 

    If $a=2$, the subtraction is $\frac32$ and the net change is
    $(1+\frac32)-\frac32=1$.  

    
    If $a>2$, subtraction is
    \[
    1 + \frac32\left(1 + \cdots + 2^{a-3}\right) = 1 + \frac32\left(
      2^{a-2}-1 \right)
    \]

    Net change:
    \[
    \left(1 + \frac32\cdot 2^{a-2} \right) - 
    \left[ 1 + \frac32\left(
      2^{a-2}-1 \right) \right] = \frac32 > 1. 
    \]

  \item \emph{One of the zeros is at $-a$ and the other is at $b$ with
      $a,b>1$.}

    Addition to weight: $\frac32\cdot2^{a-2} + \frac32\cdot2^{b-2} =
    \frac32\left(2^{a-2}+2^{b-2}\right)$. 

    \begin{itemize}
    \item If $a=b=2$: Addition becomes $\frac32 + \frac32$.
      Subtraction is $1+1=2$. Net is $3-2=1$.


    \item If $a>2$ and $b=2$: Addition becomes
      $\frac32\cdot2^{a-2}+\frac32$. 

      On the left, subtraction is 
      \[
      1 + \frac32\left(1+ \cdots + 2^{a-3}\right) = 
      1 +\frac32\left(2^{a-2}-1\right)
      \]
      On the right, the subtraction is $1$. Complete subtraction is
      \[
      2+\frac32\left(2^{a-2}-1\right) = \frac12 + \frac32\cdot2^{a-2}.
      \]
      
      Therefore net change to weight is 
      \[
      \left( \frac32\cdot2^{a-2}+\frac32 \right) - 
      \left( \frac12 + \frac32\cdot2^{a-2} \right) = 1.
      \]

    \item If $a=2$ and $b>2$: Same as previous case.

    \item If $a>2$ and $b>2$:


    Subtraction from the weight. We have $1+1$ in the middle.
    On the left, the weights in positions $2$ through $a-1$ are
    $\frac32, \frac32\cdot2, \ldots, \frac32 \cdot2^{a-3}$ and
    likewise on the right. Total subtraction is
    \[
    \frac32\left(1+2+\cdots+2^{a-3}\right) + 1 + 1 +
    \frac32\left(1+2+\cdots+2^{b-3}\right)
    \]
    which equals
    \[
    \frac32\left((2^{a-2}-1)  + (2^{b-2}-1)\right) + 2 
    =
    \frac32\left(2^{a-2} + 2^{b-2}\right) -1 .
    \]
    Therefore, net change to weight is
    \[
    \frac32\left(2^{a-2}+2^{b-2}\right) - \left[
      \frac32\left(2^{a-2} + 2^{b-2}\right) -1 
    \right] = 1.
    \]
    \end{itemize}
    

  \end{itemize}

  In all cases, an H or an I transform applied to a word increases the
  weight of the word by at least $1$.
\end{proof}


\begin{prop}
  In $W_b$ we have
  \[
  \weight{111\ldots1} =  
  \begin{cases}
    3\cdot 2^{b/2-1}-1 & \text{if $b$ is even and} \\
    2^{(b+1)/2}-1 & \text{if $b$ is odd}.
  \end{cases}
  \]
\end{prop}

\begin{proof}
  If $b$ is even, let $k=b/2$. The weight of the all 1s word is
  \[
  \left(3\cdot 2^{k-2} + 3\cdot 2^{k-3} + \cdots + \frac32\right) + 1 
  + 
  1 + \left( \frac32 + \cdots + 3\cdot 2^{k-3} + 3\cdot 2^{k-2}\right)
  \]
  which equals 
  \[
  \begin{aligned}
    2\cdot3\cdot(2^{k-2} + \cdots +2+ 1 + \tfrac12) + 2    
    &= 3 \cdot(2^{k-1} + \cdots + 1) + 2 \\
    &= 3 \cdot (2^{k} -1) + 2 \\
    &= 3 \cdot 2^{b/2} - 1.
  \end{aligned}
  \]

  If $b$ is odd, let $b=2k+1$. The weight of the all 1s word is
  \[
  \left(2^{k-1}+2^{k-2}+\cdots+1\right) + 1 + 
  \left(1 + 2 + \cdots + 2^{k-2}+2^{k-1}\right)
  \]
  which equals $2\cdot(2^k-1) + 1 = 2^{k+1}-1 = 2^{(b+1)/2}-1$.
\end{proof}

\begin{proof}[Proof of Theorem~\ref{thm:main}]
  For a positive integer $b$ let
  \[
  f(b) =
  \begin{cases}
    3\cdot 2^{b/2-1} & \text{if $b$ is even and} \\
    2^{(b+1)/2} & \text{if $b$ is odd}.
  \end{cases}
  \]

  By Proposition~\ref{prop:c-lower}, $c(b)\ge f(b)$. 

  The cardinaltiy of a maximum chain in $W_b$ cannot be larger $f(b)$
  since the lowest possible weight of a word is $0$ (for the all zeros
  word) and the maximum possible weight is $f(b)-1$ (for the all ones
  word) and for each pair of sequential elements $v \prec w$ in any
  chain the $\weight{w}\ge\weight{v}+1$. 
\end{proof}


\end{document}
